\chapter{Coarse Maps I: Aggregation and Kron Closure}
\label{ch:gaussian-coarsegraining}

Which coarse maps of the Gaussian interaction family are exact, and when does
the family close under them? The core distinguishes Markov pushforward, energy
precomposition, and restriction of a recognition family
(\Cref{sec:cg-three-operations}); this chapter realizes the first two on the
interaction family of \Cref{def:gauss-interaction-family}. Hard aggregation
precomposes the energy with the identification of each cluster to one
coordinate. For graph-exponential energies it closes with a linear parameter
map (\Cref{thm:cg-graph-exponential-closure}); on the Gaussian family it sums
cut weights and self terms and annihilates internal edges
(\Cref{prop:cg-gaussian-aggregation}), and the coarse law is the fine law
conditioned on zero within-cluster disagreement
(\Cref{prop:cg-aggregation-conditioning}), which is also a precision-dependent
linear pushforward (\Cref{prop:ig-galerkin-pushforward}). Kron reduction
marginalizes hidden agents. It is exact on the complete Gaussian action
(\Cref{prop:cg-normalized-schur-action}) but leaves the matrix-weighted
interaction family in general (\Cref{prop:cg-kron-leaves-family}), and
\Cref{tab:cg-kron-regimes} collects the regimes in which it closes, among them
the graph criterion of \Cref{thm:cg-kron-boundary-classification}, its
prior-carrying form \Cref{cor:cg-kron-anchored-agents}, and the nonflat series
closure of \Cref{thm:cg-nonflat-series-kron}. These results are the Gaussian
instance of the coarse channels of \Cref{ch:coarsegraining} and of the network
coarse-graining of \Cref{ch:agent-network-rg}. Nonflat internal links,
relaxation, and quotient volumes continue in \Cref{ch:gaussian-nonflat}, the
operator-level aggregation semigroup is \Cref{prop:grg-aggregation-semigroup},
and reciprocal pairs are treated in \Cref{ch:obstructions}.

Symbols without a channel index refer to the belief channel, of fiber
dimension \(K\); a model-channel realization, of dimension \(d_m\), uses its own
self terms, weights, and links, and no equality of the two channels is implied.
Coarse families admit newly generated positive-semidefinite self terms; the
family that forbids them is the no-new-anchor pair subfamily of
\Cref{def:cg-no-new-anchor-pair}. A law has density proportional to
\(e^{-\mathsf E}\) with respect to its stated reference, as in
\eqref{eq:cg-energy-precomposition}, and every quadratic energy carries the
factor one half. The chapter uses these local symbols: \(\mathsf h\) for a
Gaussian source (information vector), \(\mathfrak c\) for a scalar offset,
subscripts \(\mathrm r\) and \(\mathrm h\) for retained and hidden (eliminated)
coordinates, \(V_{\mathrm r}\) and \(V_{\mathrm h}\) for retained and hidden
vertices, \(M\) for the pivot of an elimination, \(H\) for the frame of a
congruence-diagonal cone \(\mathcal W_H\), and \(t_1,t_2\) for positive test
scales.

This chapter does not claim that a separately fitted coarse coupling, or a
learned or soft pooling, realizes an exact coarse map; those are different
approximate models. Nor does it claim a canonical decomposition of a coarse
operator into links and self terms, or monotonicity of an information
functional along aggregation when the precision varies. Unrestricted Kron
closure is refuted, and balanced cut links admit no continuous extension to all singular cuts
(\Cref{open:cg-cut-compression}). The open items are monotonicity along aggregation with
varying precision (\Cref{open:grg-galerkin-monotonicity}) and Kron
classification beyond finite descriptions
(\Cref{open:cg-kron-classification-beyond-finite}). \status{NOT-CLAIMED}

\section{Graph-exponential energy closure}
\label{sec:cg-graph-exponential}

Hard identification of a cluster is precomposition of the energy with the map
that copies one coarse coordinate to every member. For energies built from node
and pair statistics this precomposition stays in the same form.

\definitionheading{Graph-exponential energy}{def:cg-graph-exponential-energy}
Let \(V\) be finite, let \(\mathsf Z_0\) be standard Borel with a
\(\sigma\)-finite reference \(\nu_0\), and let
\(\mathsf T_0:\mathsf Z_0\to\mathsf V_0\) and \(\mathsf T_1:\mathsf Z_0^2\to\mathsf V_1\) be
measurable statistics with values in real vector spaces, \(\mathsf T_1\) symmetric.
For node parameters \(\theta_i\) and pair parameters
\(\theta_{ij}=\theta_{ji}\), paired with \(\mathsf V_0\) and \(\mathsf V_1\),
the graph-exponential energy on \(\mathsf Z_0^V\) is
\begin{equation}
\mathsf E_\theta(z)
=-\sum_{i\in V}\langle\theta_i,\mathsf T_0(z_i)\rangle
-\sum_{\{i,j\}}\langle\theta_{ij},\mathsf T_1(z_i,z_j)\rangle .
\label{eq:cg-graph-energy}
\end{equation}
It defines a law, with density \(e^{-\mathsf E_\theta}/Z_\theta\) relative to
\(\nu_0^{\otimes V}\), exactly when
\(Z_\theta=\int e^{-\mathsf E_\theta}d\nu_0^{\otimes V}\) is finite. For a
partition \(\mathcal P\) of \(V\), with \(I(i)\) the cluster of \(i\), the
identification map is
\(\iota_{\mathcal P}:\mathsf Z_0^{\mathcal P}\to\mathsf Z_0^V\),
\((\iota_{\mathcal P}\bar z)_i=\bar z_{I(i)}\).

\theoremheading{Graph-exponential trace closure}{thm:cg-graph-exponential-closure}
Suppose the diagonal restriction of the pair statistic is affine in the node
statistic,
\begin{equation}
\mathsf T_1(z,z)=\mathsf D \mathsf T_0(z)+\mathsf T_1^\circ,
\label{eq:cg-diagonal-affinity}
\end{equation}
for a linear map \(\mathsf D:\mathsf V_0\to\mathsf V_1\) and a constant
\(\mathsf T_1^\circ\); this excludes pair statistics whose diagonal the node statistic
cannot represent. Then for every partition \(\mathcal P\),
\begin{equation}
\mathsf E_\theta(\iota_{\mathcal P}\bar z)
=-\sum_I\langle\theta_I,\mathsf T_0(\bar z_I)\rangle
-\sum_{I<J}\langle\theta_{IJ},\mathsf T_1(\bar z_I,\bar z_J)\rangle
+\mathfrak c_{\mathcal P}(\theta),
\label{eq:cg-coarse-energy}
\end{equation}
where
\begin{equation}
\theta_I=\sum_{i\in I}\theta_i
+\mathsf D^\ast\sum_{\{i,j\}\subseteq I}\theta_{ij},
\qquad
\theta_{IJ}=\sum_{\substack{i\in I\\j\in J}}\theta_{ij},
\qquad
\mathfrak c_{\mathcal P}(\theta)
=-\Bigl\langle\sum_I\sum_{\{i,j\}\subseteq I}\theta_{ij},\mathsf T_1^\circ\Bigr\rangle .
\label{eq:cg-coarse-parameters}
\end{equation}
The parameter map \(\theta\mapsto(\theta_I,\theta_{IJ})\) is linear. A coarse
law, with density proportional to \(e^{-\mathsf E_\theta\circ\iota_{\mathcal P}}\)
relative to \(\nu_0^{\otimes\mathcal P}\), exists if and only if
\begin{equation}
\bar Z_{\mathcal P}(\theta)
=\int e^{-\mathsf E_\theta(\iota_{\mathcal P}\bar z)}
\nu_0^{\otimes\mathcal P}(d\bar z)<\infty .
\label{eq:cg-coarse-normalizer}
\end{equation}
If \(\mathsf T_1(z,z)\) is constant, \(\mathsf D=0\) is admissible, internal edges
contribute only to \(\mathfrak c_{\mathcal P}\), and node parameters add. When
\(\mathsf Z_0\) is a group and \(\mathsf T_1\) is translation invariant, the diagonal
is constant; the converse fails, as \(\mathsf T_1(z,z')=(z-z')^2zz'\) on \(\R\) shows.
If no nonzero functional of \(\mathsf T_0\) is \(\nu_0\)-almost everywhere constant,
the map \(\mathsf D\) in \eqref{eq:cg-diagonal-affinity} is unique.
\status{ESTABLISHED}

\noindent\emph{Proof.}
Substitute the identified variables into \eqref{eq:cg-graph-energy}. A cut edge
\(\{i,j\}\) with \(i\in I\), \(j\in J\) contributes
\(-\langle\theta_{ij},\mathsf T_1(\bar z_I,\bar z_J)\rangle\), and these sum to the
\(\theta_{IJ}\) term. An internal edge contributes, by
\eqref{eq:cg-diagonal-affinity},
\(-\langle\mathsf D^\ast\theta_{ij},\mathsf T_0(\bar z_I)\rangle
-\langle\theta_{ij},\mathsf T_1^\circ\rangle\); the first part joins \(\theta_I\) and
the second joins \(\mathfrak c_{\mathcal P}\). This proves
\eqref{eq:cg-coarse-energy} and \eqref{eq:cg-coarse-parameters}, and the
parameter map is a sum of linear maps. The normalizer clause is the definition
of a law on the stated reference; the algebra does not decide finiteness. A
translation-invariant \(\mathsf T_1\) has \(\mathsf T_1(z,z)=\mathsf T_1(0,0)\). The displayed
\(\mathsf T_1\) vanishes on the diagonal but changes under \(z\mapsto z+a\). If two
maps \(\mathsf D,\mathsf D'\) satisfy \eqref{eq:cg-diagonal-affinity}, then
\((\mathsf D-\mathsf D')\mathsf T_0\) is almost everywhere constant, so
\(\langle(\mathsf D-\mathsf D')^\ast\varphi,\mathsf T_0\rangle\) is almost everywhere
constant for every functional \(\varphi\), and the hypothesis forces
\(\mathsf D=\mathsf D'\). \(\square\)

The theorem is energy closure, not closure of normalized laws. The reference
\(\nu_0^{\otimes\mathcal P}\), the normalizer, and any gauge action are further
data: a compatible normalized family would satisfy
\(\mathbb P_\theta C_{\mathcal P}=\bar{\mathbb P}_{\theta_{\mathcal P}}\) for a
declared channel \(C_{\mathcal P}\) and the parameter map above, and energy
precomposition does not supply that channel (\Cref{sec:cg-three-operations}).
For the Gaussian family the channel exists and is identified in
\Cref{prop:ig-galerkin-pushforward}.

\section{Hard aggregation of the interaction precision}
\label{sec:cg-gaussian-aggregation}

In trivialized coordinates \(z_i\) the interaction precision of
\Cref{def:gauss-interaction-family} and its energy are
\begin{equation}
\Lambda_{ii}=A_i+\sum_{j\neq i}W_{ij},\qquad
\Lambda_{ij}=-W_{ij}\quad(i\neq j),
\qquad
A_i\succeq0,\quad W_{ij}=W_{ji}\succeq0,
\label{eq:cg-interaction-family}
\end{equation}
\begin{equation}
\mathsf E(z)=\frac12z^\top\Lambda z
=\frac12\sum_i z_i^\top A_i z_i
+\frac12\sum_{i<j}(z_i-z_j)^\top W_{ij}(z_i-z_j),
\label{eq:cg-trivialized-energy}
\end{equation}
the second identity being \Cref{prop:gauss-interaction-energy-kernel} with the
factor one half. The form is a declared subfamily of the directed
linear-Gaussian models (\Cref{prop:gauss-global-condition-witnesses,prop:gauss-law-level-witnesses}).

\propositionheading{Exact Gaussian aggregation}{prop:cg-gaussian-aggregation}
For a hard partition \(\mathcal P\) of the \(N\) agents, let
\(\widehat S\in\{0,1\}^{N\times|\mathcal P|}\) with
\(\widehat S_{iI}=\mathbf 1_{\{i\in I\}}\), put
\(S=\widehat S\otimes I_K\) and \(\Lambda_{\mathrm c}=S^\top\Lambda S\), and let
\(n_I=|I|\). The columns of \(S\) are orthogonal but not normalized,
\begin{equation}
S^\top S=\operatorname{diag}(n_I)\otimes I_K .
\label{eq:cg-aggregation-matrix}
\end{equation}
For every \(\Lambda\) of the form \eqref{eq:cg-interaction-family},
\begin{equation}
(\Lambda_{\mathrm c})_{IJ}
=-\sum_{\substack{i\in I\\j\in J}}W_{ij}\quad(I\neq J),
\qquad
(\Lambda_{\mathrm c})_{II}
=\sum_{i\in I}A_i
+\sum_{\substack{i\in I\\j\notin I}}W_{ij},
\label{eq:cg-coarse-blocks}
\end{equation}
so \(\Lambda_{\mathrm c}\succeq0\) is again an interaction precision, with
\(A_I=\sum_{i\in I}A_i\) and \(W_{IJ}=\sum_{i\in I,j\in J}W_{ij}\); every
internal edge is annihilated, and a source \(\mathsf h\) becomes
\(S^\top\mathsf h\). \status{ESTABLISHED}

\noindent\emph{Proof.}
\Cref{eq:cg-aggregation-matrix} holds because each agent lies in exactly one
cluster. The Gaussian energy with source,
\(\mathsf E(z)-\mathsf h^\top z\), is the graph-exponential energy
\eqref{eq:cg-graph-energy} with \(\mathsf T_0(z)=(z,-\tfrac12zz^\top)\),
\(\theta_i=(\mathsf h_i,A_i)\), \(\mathsf T_1(z,z')=-\tfrac12(z-z')(z-z')^\top\), and
\(\theta_{ij}=W_{ij}\), paired by \(\langle(\mathsf h,A),(z,Z)\rangle
=\mathsf h^\top z+\Tr(AZ)\). The pair statistic vanishes on the diagonal, so
\Cref{thm:cg-graph-exponential-closure} applies with \(\mathsf D=0\) and
\(\mathsf T_1^\circ=0\) and gives \(A_I=\sum_{i\in I}A_i\),
\(\mathsf h_I=\sum_{i\in I}\mathsf h_i=(S^\top\mathsf h)_I\), and
\(W_{IJ}=\sum W_{ij}\), with no internal-edge term. Reading the blocks of the
coarse energy through \Cref{prop:gauss-interaction-energy-kernel} gives
\eqref{eq:cg-coarse-blocks}, and \(\Lambda_{\mathrm c}\succeq0\) by
congruence. \(\square\)

This is the Galerkin coarse operator of unsmoothed aggregation algebraic
multigrid, whose prolongator is the Boolean aggregate-membership matrix
\citep[Eqs.~(2.1)--(2.2)]{Notay2010Aggregation}; for graph Laplacians the same
piecewise-constant interpolation returns a graph Laplacian
\citep[Sec.~2]{LivneBrandt2012LAMG}. Smoothed aggregation smooths this
tentative prolongator and builds it from declared zero-energy modes, which
covers systems with several unknowns per node \citep{VanekMandelBrezina1996};
\(S=\widehat S\otimes I_K\) is the piecewise-constant prolongator for \(K\)
unknowns per node. The general Galerkin construction is
\Cref{sec:cg-galerkin}.

\propositionheading{Aggregation is conditioning on zero internal disagreement}{prop:cg-aggregation-conditioning}
Let \(\Lambda\succ0\) on \(\R^n\), let \(\mathsf h\in\R^n\), and let
\(p=\mathcal N(\Lambda^{-1}\mathsf h,\Lambda^{-1})\). Let
\(B\in\R^{n\times r}\) have full column rank, let the columns of \(B_\perp\)
form an orthonormal basis of \(\operatorname{range}(B)^\perp\), and put
\(w=B_\perp^\top z\). The conditional law of \(z\sim p\) given \(w\) has a
unique version that is weakly continuous in \(w\). At \(w=0\) it is supported
on \(\operatorname{range}(B)\), and in the coordinate \(y\) with \(z=By\) it is
\begin{equation}
p(dy\mid B_\perp^\top z=0)
=\mathcal N\bigl((B^\top\Lambda B)^{-1}B^\top\mathsf h,(B^\top\Lambda B)^{-1}\bigr),
\label{eq:cg-aggregation-conditional}
\end{equation}
the Gaussian with natural parameters \((B^\top\mathsf h,B^\top\Lambda B)\). The
same law results from conditioning, in the same continuous sense, on
\(\psi z=0\) for any linear map \(\psi\) with \(\ker\psi=\operatorname{range}(B)\).
The linearity is material: when \(r<n\), for a smooth \(\chi\) on \(\R^r\) with
\(\chi(y)\geq c>0\) for every \(y\), conditioning on
the statistic \(\chi(B^+z)B_\perp^\top z\), which has the same zero set,
multiplies the density of \eqref{eq:cg-aggregation-conditional} by
\(\chi(y)^{-(n-r)}\) before normalization, and so changes the law whenever
\(\chi\) is not constant. For \(B=S\) the event \(B_\perp^\top z=0\) says
\(z_i=z_j\) whenever \(i\) and \(j\) share a cluster, so
\((S^\top\mathsf h,\Lambda_{\mathrm c})\) are the natural parameters of the fine
law conditioned on zero within-cluster disagreement. When \(r<n\) this
conditional law is singular with respect to \(p\). \status{ESTABLISHED}

\noindent\emph{Proof.}
The map \((y,w)\mapsto By+B_\perp w\) is a linear bijection, so in the
coordinates \((y,w)\) the law \(p\) has a density proportional to
\(\exp[-\frac12(By+B_\perp w)^\top\Lambda(By+B_\perp w)
+\mathsf h^\top(By+B_\perp w)]\), with a constant Jacobian. Since
\(B_\perp^\top B=0\) and \(B_\perp^\top B_\perp=I\), the statistic is
\(B_\perp^\top z=w\). The joint density is continuous and positive, so dividing
by the \(w\)-marginal density gives a conditional density jointly continuous in
\((y,w)\); it is Gaussian with mean affine in \(w\) and constant covariance,
hence weakly continuous, and two versions agree for almost every \(w\), hence
for every \(w\) by continuity. At \(w=0\) the exponent is
\(-\frac12y^\top B^\top\Lambda By+(B^\top\mathsf h)^\top y\), which gives
\eqref{eq:cg-aggregation-conditional}. If \(\ker\psi=\operatorname{range}(B)\),
then \(\psi B=0\) gives \(\psi=\psi B_\perp B_\perp^\top\) with \(\psi B_\perp\)
injective, so \(\psi z=(\psi B_\perp)w\) and \(w=(\psi B_\perp)^+\psi z\); the two
statistics generate the same \(\sigma\)-algebra, and their continuous versions
correspond under \(w\mapsto\psi B_\perp w\), which fixes \(0\). For the
nonlinear statistic, \(B^+z=y\), and in the coordinates \((y,u)\) with
\(u=\chi(y)w\) the density of \(p\) acquires the Jacobian factor
\(\chi(y)^{-(n-r)}\). The lower bound on \(\chi\) and positive definiteness of
\(\Lambda\) bound these joint densities near \(u=0\) by an integrable Gaussian
function of \(y\). Dominated convergence therefore gives a positive finite marginal
density and a weakly continuous conditional version at zero, with the stated density.
For \(r<n\) this is a different law exactly when \(\chi\) is not constant.
Without an integrability condition the nonlinear statement fails: for independent standard
normal \((Y,W)\) and \(\chi(y)=e^{-y^2}\), the candidate density at \(u=0\) is
proportional to \(e^{y^2/2}\), which is not integrable. For \(B=S\), \(\operatorname{range}(S)\) is the set of configurations
constant on each cluster. Finally \(\operatorname{range}(B)\) is Lebesgue null
when \(r<n\), while \(p\) has a positive density. \(\square\)

\propositionheading{Galerkin aggregation is a precision-dependent pushforward}{prop:ig-galerkin-pushforward}
Let \(\Lambda\succ0\) on \(\R^n\), let \(S\in\R^{n\times r}\) be injective,
and put
\begin{equation}
\Lambda_{\mathrm c}=S^\top\Lambda S,
\qquad
P_\Lambda=\Lambda_{\mathrm c}^{-1}S^\top\Lambda .
\label{eq:ig-galerkin-projector}
\end{equation}
(i) The map \(y\mapsto P_\Lambda y\) pushes
\(\mathcal N(\Lambda^{-1}\mathsf h,\Lambda^{-1})\) to
\(\mathcal N(\Lambda_{\mathrm c}^{-1}S^\top\mathsf h,\Lambda_{\mathrm c}^{-1})\),
the Gaussian with natural parameters \((S^\top\mathsf h,S^\top\Lambda S)\)
obtained by precomposing the energy with \(y=Sc\).
(ii) On the fixed-precision mean fiber
\(\{\mathcal N(\mu,\Lambda^{-1}):\mu\in\R^n\}\), \(P_\Lambda\) is one
parameter-independent deterministic kernel, and the Fisher metric of the pushed
family satisfies
\(P_\Lambda^\top\Lambda_{\mathrm c}P_\Lambda
=\Lambda S\Lambda_{\mathrm c}^{-1}S^\top\Lambda\preceq\Lambda\), with equality
exactly on tangent directions in \(\operatorname{range}(S)\).
(iii) For nested aggregations \(S=S_1S_2\) with both factors injective,
\(P^{(S_2)}_{S_1^\top\Lambda S_1}P^{(S_1)}_\Lambda=P^{(S_1S_2)}_\Lambda\), where
the superscript names the aggregation matrix.
(iv) Let the joint law of \((o,Y)\) given the structural data be Gaussian with
precision \(\widehat\Lambda\) and information vector \(\widehat{\mathsf h}\), both
functions of \((\theta,X)\) alone, so that by
\Cref{prop:restrict-conditional-precision-observation-free} the conditional
precision of \(Y\) given \(o\) is \(\widehat\Lambda_{YY}\). Then
\(\operatorname{id}_O\otimes P_{\widehat\Lambda_{YY}}\) is a recognition-independent
channel determined by \((\theta,X)\). It leaves the observation marginal, hence
the evidence, unchanged, and it carries the conditional law of \(Y\) given \(o\)
to the Gaussian with natural parameters
\((S^\top(\widehat{\mathsf h}_Y-\widehat\Lambda_{Yo}o),S^\top\widehat\Lambda_{YY}S)\), which is the
conditional law of \(c\) given \(o\) under the joint precomposed with \(y=Sc\);
that precomposed joint generally has a different observation marginal.
(v) Let \(\mathfrak L\) be a set of positive-definite \(n\times n\) matrices. A
single Markov kernel \(\kappa\) with
\(\mathcal N(\mu,\Lambda^{-1})\kappa=\mathcal N(P_\Lambda\mu,\Lambda_{\mathrm c}^{-1})\)
for every \(\mu\in\R^n\) and every \(\Lambda\in\mathfrak L\) exists if and only
if the subspace \(\operatorname{range}(\Lambda S)\) is the same for every
\(\Lambda\in\mathfrak L\); then \(P_\Lambda\) is constant on \(\mathfrak L\) and
\(\kappa(y,\cdot)=\delta_{P_\Lambda y}\) for Lebesgue-almost every \(y\). For a
hard partition with a cluster of at least two agents, the interaction precision
with every \(A_i=I_K\) and no edges, and the one obtained from it by replacing
one self term in that cluster by \(3I_K\), give different subspaces, so no
parameter-independent kernel realizes aggregation on the interior of the
interaction cone. \status{ESTABLISHED}

\noindent\emph{Proof.}
(i) \(P_\Lambda\Lambda^{-1}P_\Lambda^\top
=\Lambda_{\mathrm c}^{-1}S^\top\Lambda\Lambda^{-1}\Lambda S\Lambda_{\mathrm c}^{-1}
=\Lambda_{\mathrm c}^{-1}\) and
\(P_\Lambda\Lambda^{-1}\mathsf h=\Lambda_{\mathrm c}^{-1}S^\top\mathsf h\), and a
linear image of a Gaussian is Gaussian. (ii) The pushed family
\(\mu\mapsto\mathcal N(P_\Lambda\mu,\Lambda_{\mathrm c}^{-1})\) has Fisher matrix
\(P_\Lambda^\top\Lambda_{\mathrm c}P_\Lambda\), and
\(\Lambda-\Lambda S\Lambda_{\mathrm c}^{-1}S^\top\Lambda
=\Lambda^{1/2}(I-\Pi)\Lambda^{1/2}\) with \(\Pi\) the orthogonal projection onto
\(\operatorname{range}(\Lambda^{1/2}S)\), whose null directions are exactly
\(\operatorname{range}(S)\). (iii) Substitute \(\Lambda_1=S_1^\top\Lambda S_1\)
into \((S_2^\top\Lambda_1S_2)^{-1}S_2^\top\Lambda_1\Lambda_1^{-1}S_1^\top\Lambda\).
(iv) The channel is the identity on \(o\) and depends only on
\(\widehat\Lambda_{YY}\), a function of \((\theta,X)\). The conditional law is
\(\mathcal N(\widehat\Lambda_{YY}^{-1}(\widehat{\mathsf h}_Y-\widehat\Lambda_{Yo}o),\widehat\Lambda_{YY}^{-1})\),
so (i) with source \(\widehat{\mathsf h}_Y-\widehat\Lambda_{Yo}o\) gives the stated image, and
reading the precomposed joint density at fixed \(o\) as a function of \(c\) gives
the same natural parameters. (v) Suppose \(\kappa\) exists and fix
\(\Lambda_0\in\mathfrak L\). For a bounded continuous \(f\) on \(\R^r\) put
\(F(y)=\int f(u)\kappa(y,du)\) and \(F_0=f\circ P_{\Lambda_0}\). By the
hypothesis on \(\kappa\) and by (i), the Gaussian averages of \(F\) and of
\(F_0\) against \(\mathcal N(\mu,\Lambda_0^{-1})\) agree for every \(\mu\), so
the bounded function \(F-F_0\) has zero convolution with the
\(\mathcal N(0,\Lambda_0^{-1})\) density. Its Fourier transform is a tempered
distribution annihilated by multiplication with the nowhere-vanishing Gaussian
transform, hence zero, and \(F=F_0\) almost everywhere. Running \(f\) over
\(u\mapsto e^{i\omega^\top u}\) with rational \(\omega\) and using continuity of
characteristic functions gives \(\kappa(y,\cdot)=\delta_{P_{\Lambda_0}y}\) for
almost every \(y\). Gaussian laws are absolutely continuous, so the hypothesis
at any \(\Lambda\in\mathfrak L\) requires \(P_{\Lambda_0}\mu=P_\Lambda\mu\) for
every \(\mu\). Both maps are left inverses of \(S\), with kernels
\(\operatorname{range}(\Lambda_0S)^\perp\) and
\(\operatorname{range}(\Lambda S)^\perp\), and a left inverse of \(S\) is
determined by its kernel, so they agree exactly when
\(\operatorname{range}(\Lambda S)=\operatorname{range}(\Lambda_0S)\). Conversely,
if \(\Lambda S=\Lambda_0SG\) with \(G\) invertible, then
\(P_\Lambda=P_{\Lambda_0}\) and \(\Lambda^{-1}\Lambda_0S=SG^{-1}\) give
\(P_{\Lambda_0}\Lambda^{-1}P_{\Lambda_0}^\top=G^{-1}(S^\top\Lambda_0S)^{-1}
=\Lambda_{\mathrm c}^{-1}\), so \(\delta_{P_{\Lambda_0}y}\) serves every
\(\Lambda\in\mathfrak L\). For the final clause, the first precision is the
identity, with \(\operatorname{range}(\Lambda S)=\operatorname{range}(S)\), the
vectors constant on every cluster. For the second, apply \(\Lambda S\) to the
coarse vector equal to \(v\neq0\) on the cluster and zero elsewhere: the result
equals \(3v\) at one agent of the cluster and \(v\) at another, so it is not
constant on the cluster and \(\operatorname{range}(\Lambda S)\neq
\operatorname{range}(S)\). \(\square\)

By part (i), the conditional law of \Cref{prop:cg-aggregation-conditioning} is
the image of \(p\) under \(P_\Lambda\), which depends on \(\Lambda\). With
\(B=\iota_{\mathcal P}\), the injection of \Cref{sec:cg-partial},
\Cref{prop:cg-aggregation-conditioning} identifies
\(\iota_{\mathcal P}^\top\Lambda\iota_{\mathcal P}\) as the precision of the fine
law conditioned on exact internal parallel transport, the law-level reading of
\eqref{eq:cg-partial-properness}. The Schur term of
\eqref{eq:ig-restriction-marginal} is the data-processing defect of the mean
submodel under \(y\mapsto B^\top y\), and part (ii) exhibits the corresponding
defect \(\Lambda-\Lambda S\Lambda_{\mathrm c}^{-1}S^\top\Lambda\) for
\(P_\Lambda\). When an inference history moves within one fixed-precision mean
fiber, the Galerkin map acts pointwise as the pushforward of the kernel
\(P_\Lambda\), so the pointwise-Markov hypothesis of
\Cref{thm:hist-pointwise-contraction-lift} is met; part (v) records what fails
when the precision moves.

Part (ii) of \Cref{prop:ig-galerkin-pushforward} gives Fisher contraction on
each fixed-precision fiber, and part (iii) gives it at every finite depth of a
nested aggregation. Families on which \(\operatorname{range}(\Lambda S)\) varies,
the singular boundary member of the biadditive family, and unbounded depth are
\Cref{open:grg-galerkin-monotonicity}.

\section{Kron reduction: coefficient families and graph boundaries}
\label{sec:cg-kron}

Kron reduction marginalizes hidden agents, and it is the dual of aggregation.
For a coordinate split \(z=(z_{\mathrm r},z_{\mathrm h})\),
\Cref{prop:cg-aggregation-conditioning} with
\(\operatorname{range}(B)=\{z_{\mathrm h}=0\}\) gives the conditional precision
\(\Lambda_{\mathrm{rr}}\), the restriction of the precision. Integrating out
\(z_{\mathrm h}\) keeps instead the restriction
\((\Lambda^{-1})_{\mathrm{rr}}\) of the covariance, whose inverse is the Schur
complement
\(\Lambda_{\mathrm{rr}}-\Lambda_{\mathrm{rh}}\Lambda_{\mathrm{hh}}^{-1}\Lambda_{\mathrm{hr}}\)
(\Cref{lem:gauss-schur}), and \Cref{prop:restrict-marginal-vs-restricted-precision}
orders the two precisions. Aggregation stays in the interaction family;
marginalization does not in general (\Cref{prop:cg-kron-leaves-family}). The
regimes in which it does close carry different quantifiers: which coefficient
families pass tests at one fixed hidden anchor, which fixed graphs close for
every independent positive-definite assignment, and which nonflat graphs close
while keeping their link marks. In every regime a Schur reduction is
integration over the hidden coordinates with the source and scalar offset
carried along (\Cref{prop:cg-normalized-schur-action}), and a component whose
elimination generates a positive-semidefinite self term stays in the coarse
family. \Cref{tab:cg-kron-regimes} collects the verdicts. For scalar weights
the closure of loopy Laplacians under Kron reduction is
\citet[Lemma~II.1]{DorflerBullo2013Kron}, and
\citet[Sec.~4.2, Thm.~4.3]{HansenGhrist2019} treat Kron reduction for cellular
sheaves; the positive-pair obstruction below is compared with their results
where it is stated.

\begin{sciencetable}{Closure regimes of Kron reduction in the Gaussian interaction family. ``Generated self terms'' records whether elimination produces new positive-semidefinite self terms, which the coarse family admits.}
\label{tab:cg-kron-regimes}
\begin{tabularx}{\textwidth}{>{\raggedright\arraybackslash}X>{\raggedright\arraybackslash}p{0.24\textwidth}>{\raggedright\arraybackslash}p{0.15\textwidth}>{\raggedright\arraybackslash}p{0.19\textwidth}}
\toprule
\textbf{Admissible family and quantifier} & \textbf{Verdict} & \textbf{Generated self terms} & \textbf{Statement} \\
\midrule
Independent matrix weights and anchors on an arbitrary graph & not closed; rational three-agent witness & not applicable & \Cref{prop:cg-kron-leaves-family} \\
All self terms and weights in one fixed cone \(\mathcal W_H\) & closed on every graph & yes, inside \(\mathcal W_H\) & \Cref{thm:cg-congruence-diagonal-kron} \\
Pair tests of an arbitrary family at one fixed SPD hidden anchor \(Q_0\) & closed iff the family lies in a cone \(\mathcal W_H\) with \(HH^\top=Q_0\); maximal families are these cones & yes, inside \(\mathcal W_H\) & \Cref{thm:cg-independent-kron-classification} \\
Coordinates of \(\mathcal W_H\) confined to prescribed domains & closed iff every star update stays in the domains & yes & \Cref{prop:cg-kron-scalar-domain} \\
Positive-semidefinite hidden anchor sharing its kernel with the family & quotient form of the fixed-anchor classification & yes, on the quotient & \Cref{cor:cg-kron-shared-kernel} \\
Fixed flat graph, independent SPD weights, \(K\geq2\) & closed iff every hidden component has at most two terminals, ground counted & anchor increments from components touching ground & \Cref{thm:cg-kron-boundary-classification} \\
Fixed flat graph, SPD anchor on every hidden agent, \(K\geq2\) & closed iff every hidden component touches at most one retained agent; failure generic otherwise & SPD increment at the touched agent & \Cref{cor:cg-kron-anchored-agents} \\
Nonflat links, hidden vertices of degree two with zero anchor & closed in every order, loop holonomy retained & none & \Cref{thm:cg-nonflat-series-kron} \\
Nonflat links, hidden vertices of degree two with PSD anchors & one-link-plus-anchor presentation exists, not unique & yes & \Cref{prop:cg-nonflat-series-anchored} \\
Branching (three terminals), proper SPD law & marginal outside every positive two-node presentation; rational dual certificate & not applicable & \Cref{prop:cg-kron-positive-pair-obstruction} \\
Finite semialgebraic description of fine and coarse families & decidable by quantifier elimination & encoded in the description & \Cref{thm:cg-semialgebraic-kron-classification} \\
Declared quotient metric with variable support & decidable by quantifier elimination & encoded in the description & \Cref{thm:cg-variable-support-classification} \\
\bottomrule
\end{tabularx}
\end{sciencetable}

\propositionheading{Unrestricted Kron reduction leaves the family}{prop:cg-kron-leaves-family}
Take \(N=3\), \(K=2\), \(A_1=A_2=A_3=I_2\), and
\begin{equation}
W_{12}=0,\qquad
W_{13}=\begin{pmatrix}1&0\\0&2\end{pmatrix},\qquad
W_{23}=\begin{pmatrix}2&1\\1&2\end{pmatrix}.
\label{eq:cg-kron-witness}
\end{equation}
Then the assembled \(\Lambda\) is positive definite, and eliminating agent
\(3\) manufactures
\begin{equation}
W^{\mathrm{Kron}}_{12}=-\bigl(\Lambda^{\mathrm{eff}}\bigr)_{12}
=W_{13}\Lambda_{33}^{-1}W_{23}
=\frac1{19}\begin{pmatrix}9&3\\4&14\end{pmatrix},
\label{eq:cg-kron-asymmetric}
\end{equation}
which is not symmetric and hence is not a positive-semidefinite interaction
weight. \status{ESTABLISHED}

\noindent\emph{Proof.}
\(\Lambda_{33}=\begin{psmallmatrix}4&1\\1&5\end{psmallmatrix}\) has determinant
\(19\), and direct multiplication gives \eqref{eq:cg-kron-asymmetric}. The
leading principal minors of \(\Lambda\) are \(2,6,18,48,102,233\), so
\(\Lambda\succ0\) by Sylvester's criterion. The full Schur complement is
symmetric, its two mirrored off-diagonal blocks being transposes; membership in
\eqref{eq:cg-interaction-family} requires each block to be symmetric, which
fails. \(\square\)

\lemmaheading{Scalar Kron closure}{lem:cg-scalar-kron-closure}
For \(K=1\), call a symmetric matrix with nonpositive off-diagonal entries and
nonnegative row sums a loopy Laplacian. The Schur complement of a loopy
Laplacian at any nonsingular principal block is a loopy Laplacian; one-node
elimination with pivot \(m=\Lambda_{kk}>0\) adds \(w_{ik}w_{kj}/m\geq0\) to the
weight between distinct retained neighbors \(i,j\) of \(k\).
\status{ESTABLISHED}

\noindent\emph{Proof.}
\citet[Lemma~II.1]{DorflerBullo2013Kron} prove the statement for an irreducible
loopy Laplacian and a proper retained set of at least two nodes; the following
argument covers every case. A loopy Laplacian is diagonally dominant with
nonnegative diagonal, hence positive semidefinite, so a nonsingular principal
block is positive definite; its successive pivots are therefore positive, and
the Schur complement at the block equals the composition of one-node
eliminations. Eliminating \(k\) replaces \(\Lambda_{ij}\) by
\(\Lambda_{ij}-\Lambda_{ik}\Lambda_{kj}/m\). For \(i\neq j\) the subtracted term
is nonnegative because \(\Lambda_{ik},\Lambda_{kj}\leq0\), so off-diagonal
entries stay nonpositive and \(w_{ij}\) increases by \(w_{ik}w_{kj}/m\). With
row sums \(\varrho_i=\sum_j\Lambda_{ij}\), the new row sum of a retained \(i\) is
\(\varrho_i-\Lambda_{ik}\varrho_k/m\geq\varrho_i\geq0\). Symmetry is preserved.
\(\square\)

The matrix analog is closed on a larger subfamily than common orthogonal
diagonalization.

\definitionheading{Congruence-diagonal cone}{def:cg-congruence-diagonal-cone}
For \(H\in\GL(K)\), the congruence-diagonal cone is
\begin{equation}
\mathcal W_H
=\{H\operatorname{diag}(d)H^\top:d\in[0,\infty)^K\}.
\label{eq:cg-congruence-diagonal-cone}
\end{equation}

\theoremheading{Congruence-diagonal Kron closure}{thm:cg-congruence-diagonal-kron}
If every \(A_i\) and \(W_{ij}\) lies in one cone \(\mathcal W_H\), then every
Schur complement at a nonsingular block of hidden agents is again an interaction
precision with all reduced self terms and weights in \(\mathcal W_H\). The cones
\(\mathcal W_H\) contain pairs that no common orthogonal eigenbasis
diagonalizes: for \(H=\begin{psmallmatrix}1&1\\0&1\end{psmallmatrix}\),
\(W=H\operatorname{diag}(1,2)H^\top\) and \(W'=H\operatorname{diag}(3,1)H^\top\)
satisfy \(WW'-W'W=\begin{psmallmatrix}0&-5\\5&0\end{psmallmatrix}\).
\status{ESTABLISHED}

\noindent\emph{Proof.}
Factor
\begin{equation}
\Lambda=(I_N\otimes H)\Lambda_{\mathrm d}(I_N\otimes H^\top),
\label{eq:cg-kron-congruence-factor}
\end{equation}
where, after grouping coordinates by channel, \(\Lambda_{\mathrm d}\) is a
direct sum of \(K\) scalar loopy Laplacians. For retained agents
\(V_{\mathrm r}\) and hidden agents \(V_{\mathrm h}\),
\begin{equation}
\Lambda^{\mathrm{eff}}
=(I_{V_{\mathrm r}}\otimes H)\Bigl[
(\Lambda_{\mathrm d})_{\mathrm{rr}}
-(\Lambda_{\mathrm d})_{\mathrm{rh}}(\Lambda_{\mathrm d})_{\mathrm{hh}}^{-1}
(\Lambda_{\mathrm d})_{\mathrm{hr}}
\Bigr](I_{V_{\mathrm r}}\otimes H^\top),
\label{eq:cg-kron-congruence-schur}
\end{equation}
and \((\Lambda_{\mathrm d})_{\mathrm{hh}}\) is nonsingular because
\(\Lambda_{\mathrm{hh}}\) is. The bracket eliminates channel by channel, and
each channel's Schur complement is a loopy Laplacian by
\Cref{lem:cg-scalar-kron-closure}. Reassembling the channels gives reduced self
terms and weights in \(\mathcal W_H\). The commutator is direct multiplication,
and commuting symmetric matrices are exactly those with a common orthogonal
eigenbasis. \(\square\)

\subsection{The complete Gaussian marginal action}
\label{sec:cg-kron-normalized-action}

Fix a coordinate split \(z=(z_{\mathrm r},z_{\mathrm h})\), with \(n=\dim z\)
and \(n_{\mathrm h}=\dim z_{\mathrm h}\), an SPD precision \(\Lambda\), a source
\(\mathsf h\), and a real scalar offset \(\mathfrak c\), with coordinate
Lebesgue reference measures.

\propositionheading{Normalized Schur marginal and reverse kernel}{prop:cg-normalized-schur-action}
The action, its law, and its raw mass are
\begin{equation}
\begin{aligned}
\mathsf E(z)&=\frac12z^\top\Lambda z-\mathsf h^\top z+\mathfrak c,\\
p_\Lambda(z)&=Z_\Lambda^{-1}e^{-\mathsf E(z)},\qquad
\mathbb P_\Lambda=\mathcal N(\Lambda^{-1}\mathsf h,\Lambda^{-1}),\\
Z_\Lambda&=e^{-\mathfrak c}(2\pi)^{n/2}
(\det\Lambda)^{-1/2}
\exp\Bigl(\frac12\mathsf h^\top\Lambda^{-1}\mathsf h\Bigr).
\end{aligned}
\label{eq:cg-kron-normalized-prior}
\end{equation}
With \(M=\Lambda_{\mathrm{hh}}\), the action decomposes as
\begin{equation}
\begin{aligned}
\mathsf E(x,z_{\mathrm h})
={}&\frac12
\bigl(z_{\mathrm h}-M^{-1}(\mathsf h_{\mathrm h}-\Lambda_{\mathrm{hr}}x)\bigr)^\top
M\bigl(z_{\mathrm h}-M^{-1}(\mathsf h_{\mathrm h}-\Lambda_{\mathrm{hr}}x)\bigr)\\
&+\frac12x^\top\Lambda_{\mathrm{rr}}x-\mathsf h_{\mathrm r}^\top x+\mathfrak c
-\frac12(\mathsf h_{\mathrm h}-\Lambda_{\mathrm{hr}}x)^\top
M^{-1}(\mathsf h_{\mathrm h}-\Lambda_{\mathrm{hr}}x),
\end{aligned}
\label{eq:cg-kron-completed-square}
\end{equation}
and the exact marginal action and reverse kernel are
\begin{equation}
\begin{aligned}
\Lambda^{\mathrm{eff}}
 &=\Lambda_{\mathrm{rr}}-\Lambda_{\mathrm{rh}}M^{-1}\Lambda_{\mathrm{hr}},\\
\mathsf h^{\mathrm{eff}}
 &=\mathsf h_{\mathrm r}-\Lambda_{\mathrm{rh}}M^{-1}\mathsf h_{\mathrm h},\\
\mathfrak c^{\mathrm{eff}}
 &=\mathfrak c+\frac12\log\det M-\frac{n_{\mathrm h}}{2}\log(2\pi)
                     -\frac12\mathsf h_{\mathrm h}^\top M^{-1}\mathsf h_{\mathrm h},\\
\mathbb P_\Lambda(dz_{\mathrm h}\mid z_{\mathrm r}=x)
 &=\mathcal N\bigl(M^{-1}(\mathsf h_{\mathrm h}-\Lambda_{\mathrm{hr}}x),M^{-1}\bigr)(dz_{\mathrm h}).
\end{aligned}
\label{eq:cg-kron-complete-schur}
\end{equation}
Both \(M\) and \(\Lambda^{\mathrm{eff}}\) are SPD. Integration preserves the
raw mass \(Z_\Lambda\), and the normalized marginal is the coordinate
pushforward of \(\mathbb P_\Lambda\). Successive eliminations agree with direct
elimination, sources and scalar offsets included, and the conditional
reconstruction kernels compose. For an observation
\(o=\mathsf C_{\mathrm r}z_{\mathrm r}+\mathsf C_{\mathrm h}z_{\mathrm h}+\xi\) with
\(\xi\sim\mathcal N(0,\Sigma_o)\) independent of \(z\) and \(\Sigma_o\succ0\),
\begin{equation}
o\mid z_{\mathrm r}=x\sim
\mathcal N\Bigl(
(\mathsf C_{\mathrm r}-\mathsf C_{\mathrm h}M^{-1}\Lambda_{\mathrm{hr}})x
+\mathsf C_{\mathrm h}M^{-1}\mathsf h_{\mathrm h},\
\Sigma_o+\mathsf C_{\mathrm h}M^{-1}\mathsf C_{\mathrm h}^\top
\Bigr),
\label{eq:cg-kron-induced-sensor}
\end{equation}
with SPD covariance, and mixing this likelihood with the marginal recovers the
evidence at every record, so the coordinate projection pushes the whole
observed joint. The posterior hidden conditional differs in general from the
prior conditional in \eqref{eq:cg-kron-complete-schur}, and the two agree when
\(\mathsf C_{\mathrm h}=0\). \status{ESTABLISHED}

\noindent\emph{Proof.}
The moment formulas are the standard Gaussian marginal and conditional
\citep[Secs.~2.3.1--2.3.2, Eqs.~(2.73), (2.75), (2.91)--(2.93)]{Bishop2006};
the raw mass and scalar offset are recorded here. Translating \(z\) by
\(\Lambda^{-1}\mathsf h\) and diagonalizing \(\Lambda\) orthogonally reduce
\(\int e^{-\mathsf E}\) to one-dimensional integrals
\(\int_\R e^{-\lambda u^2/2}du=(2\pi/\lambda)^{1/2}\), which gives \(Z_\Lambda\).
Expanding the square proves \eqref{eq:cg-kron-completed-square}; integrating
its first line over \(z_{\mathrm h}\) gives
\((2\pi)^{n_{\mathrm h}/2}(\det M)^{-1/2}\), and expanding its last term proves
every line of \eqref{eq:cg-kron-complete-schur}. A principal block of an SPD
matrix is SPD, and for nonzero \(x\), minimizing
\((x,z_{\mathrm h})^\top\Lambda(x,z_{\mathrm h})\) over \(z_{\mathrm h}\) gives
\(x^\top\Lambda^{\mathrm{eff}}x>0\). For a nonnegative Borel function on a
finite product of Euclidean spaces the product integral equals its iterated
integrals in any order, with infinite values allowed; the Gaussian integrals
here are finite, which proves equality of elimination orders and the
composition of reconstruction kernels formed from the corresponding marginals.
For the sensor, write \(z_{\mathrm h}\) given \(z_{\mathrm r}=x\) as its
conditional mean plus independent centered Gaussian noise of covariance
\(M^{-1}\), which gives \eqref{eq:cg-kron-induced-sensor}
\citep[Sec.~2.3.3, Eq.~(2.115)]{Bishop2006}; integrating the likelihood of the
full coordinates against the prior in the order hidden first shows that the
evidence is unchanged. When \(\mathsf C_{\mathrm h}\neq0\) the likelihood depends
on \(z_{\mathrm h}\), which changes the hidden conditional given \(o\).
\(\square\)

Because the whole observed joint is pushed, the finite-KL chain identity of
\Cref{sec:cg-kl-recovery} gives the exact free-energy difference with the same
evidence on both sides.

\subsection{Arbitrary independent families at a fixed anchor}
\label{sec:cg-independent-kron-families}

\theoremheading{Fixed-anchor classification of arbitrary families}{thm:cg-independent-kron-classification}
Fix \(Q_0=Q_0^\top\succ0\) and an arbitrary subset
\(\mathscr F\subseteq\PSD^K\); no closedness, convexity, scaling closure, or
membership of \(Q_0\) in \(\mathscr F\) is assumed. For
\(W,W'\in\mathscr F\) and \(t_1,t_2>0\), the two-edge test eliminates a hidden
agent with anchor \(Q_0\) and incident weights \(t_1W\), \(t_2W'\), with pivot
and manufactured coefficient
\begin{equation}
M=Q_0+t_1W+t_2W',\qquad
W_{W,W'}(t_1,t_2)=t_1t_2WM^{-1}W'.
\label{eq:cg-independent-pair-test}
\end{equation}
Giving both retained agents the anchor \(Q_0\) makes each test a proper
three-agent Gaussian, with energy at least \(\lambda_{\min}(Q_0)\|z\|^2/2\).
Symmetry of \(W_{W,W'}(t_1,t_2)\) for at least one pair of scales, for every
\(W,W'\in\mathscr F\), is equivalent to
\begin{equation}
WQ_0^{-1}W'=W'Q_0^{-1}W
\quad\text{for all }W,W'\in\mathscr F,
\label{eq:cg-anchor-commutation}
\end{equation}
and to the existence of one \(O\in O(K)\) with
\begin{equation}
H=Q_0^{1/2}O,\qquad Q_0=HH^\top,\qquad
\mathscr F\subseteq\mathcal W_H .
\label{eq:cg-anchored-common-cone}
\end{equation}
One \(H\) serves the whole family, even an uncountable one. Under these
conditions every pair of positive scales gives
\(W_{W,W'}(t_1,t_2)\in\mathcal W_H\), and the inclusion-maximal families with
the test property are exactly the cones \(\mathcal W_H\). In particular, if a
closed convex cone \(\mathcal W\subseteq\PSD^K\) has an SPD order unit \(M_0\)
in its relative interior, its anchors and edge coefficients range independently
over \(\mathcal W\) with arbitrary positive scaling, and every defined one-agent
elimination must manufacture minus a symmetric member of \(\mathcal W\), then
\(\mathcal W\subseteq\mathcal W_H\) for some \(H\) with \(HH^\top=M_0\), and
every maximal such cone is a cone \(\mathcal W_H\). Neither premise can be
dropped. If only equal pairs \(W=W'\) are tested, every test output is
symmetric for the family
\(\{I,\operatorname{diag}(1,2),\begin{psmallmatrix}2&1\\1&2\end{psmallmatrix}\}\),
which lies in no common cone \(\mathcal W_H\). If the weights lie in a common
cone but the anchor does not, the test can fail: for
\(W=\operatorname{diag}(1,0)\), \(W'=\operatorname{diag}(0,1)\),
\(Q_0=\begin{psmallmatrix}2&1\\1&2\end{psmallmatrix}\), and \(t_1=t_2=1\), the
manufactured coefficient is
\(\begin{psmallmatrix}0&-1/8\\0&0\end{psmallmatrix}\). \status{ESTABLISHED}

\noindent\emph{Proof.}
Whiten by \(Q_0\), writing \(\widehat W=Q_0^{-1/2}WQ_0^{-1/2}\),
\(\widehat W'=Q_0^{-1/2}W'Q_0^{-1/2}\), and
\(\widehat M=I+t_1\widehat W+t_2\widehat W'\). Substituting
\(\widehat W'=(\widehat M-I-t_1\widehat W)/t_2\) gives
\begin{equation}
\begin{aligned}
\widehat W\widehat M^{-1}\widehat W'-\widehat W'\widehat M^{-1}\widehat W
 &=t_2^{-1}(\widehat M^{-1}\widehat W-\widehat W\widehat M^{-1})
  =\widehat M^{-1}(\widehat W\widehat W'-\widehat W'\widehat W)\widehat M^{-1},\\
W_{W,W'}-W_{W,W'}^\top
 &=t_1t_2Q_0M^{-1}(WQ_0^{-1}W'-W'Q_0^{-1}W)M^{-1}Q_0.
\end{aligned}
\label{eq:cg-exact-pair-antisymmetry}
\end{equation}
All outside factors are invertible or positive scalars, so one tested pair of
positive scales already forces \eqref{eq:cg-anchor-commutation}. The whitened
family is real symmetric and pairwise commuting, and simultaneous orthogonal
diagonalization of the whole family follows by induction on dimension: if every
member is scalar any orthonormal basis works; otherwise every member preserves
the eigenspaces of one nonscalar member, which have smaller dimension, and the
induction applies to the restricted families. Only finitely many reductions are
used, so no countability premise is needed, and the empty family is trivial.
Positive semidefiniteness makes the diagonal entries nonnegative, which proves
\eqref{eq:cg-anchored-common-cone}. Conversely, if \(W=H\operatorname{diag}(d)H^\top\)
and \(W'=H\operatorname{diag}(d')H^\top\) with \(d,d'\in[0,\infty)^K\), then
\begin{equation}
W_{W,W'}(t_1,t_2)=
H\operatorname{diag}(v)H^\top,
\qquad
v=\frac{t_1t_2 d\odot d'}{\mathbf 1+t_1d+t_2d'}\ \text{coordinatewise},
\label{eq:cg-independent-diagonal-output}
\end{equation}
which proves sufficiency and positivity. A tested family therefore lies in a
full cone \(\mathcal W_H\), and a maximal family equals it. Conversely, a PSD
matrix compatible with every member of \(\mathcal W_H\) becomes, after
whitening and the same orthogonal change, a matrix commuting with every
coordinate rank-one projector, hence nonnegative diagonal, so it belongs to the
cone; this proves maximality.

For the cone statement, \(M_0\in\mathcal W\) and scaling closure admit the pair
test with \(Q_0=M_0\), incident weights \(t_1W,t_2W'\) for
\(W,W'\in\mathcal W\), and retained anchors \(M_0\); the required symmetry
gives \eqref{eq:cg-anchor-commutation} with \(Q_0=M_0\), hence
\(\mathcal W\subseteq\mathcal W_H\) with \(HH^\top=M_0\). The cone
\(\mathcal W_H\) is closed and convex, contains \(M_0=HI_KH^\top\) in its
relative interior, and is closed under elimination by
\Cref{thm:cg-congruence-diagonal-kron}, so it satisfies the same hypotheses and
every maximal cone is of this form. For the omitted premises, each
\(F(I+tF)^{-1}F\) is symmetric, while \(\mathcal W_H\ni I\) forces the common
orthogonal basis \(HD_0^{1/2}\) for the family, and
\(\operatorname{diag}(1,2)\) and \(\begin{psmallmatrix}2&1\\1&2\end{psmallmatrix}\)
do not commute; the last coefficient is direct multiplication with
\(M=\begin{psmallmatrix}3&1\\1&3\end{psmallmatrix}\). \(\square\)

Anchors that vary with a context over the base are treated in Supplement B.

\remarkheading{Stability inside a prescribed diagonal domain}{prop:cg-kron-scalar-domain}
Fix \(H\) and coordinate domains
\(\mathscr D_A,\mathscr D_W\subseteq[0,\infty)^K\) with
\(0\in\mathscr D_W\), and write \(A_i=H\operatorname{diag}(a_i)H^\top\) and
\(W_{ij}=H\operatorname{diag}(w_{ij})H^\top\). A defined elimination of agent
\(k\) has coordinatewise pivot \(m=a_k+\sum_{i\neq k}w_{ik}>0\) and updates
\begin{equation}
w'_{ij}=w_{ij}+\frac{w_{ik}\odot w_{kj}}{m},\qquad
a'_i=a_i+\frac{w_{ik}\odot a_k}{m},
\label{eq:cg-kron-scalar-updates}
\end{equation}
with products and quotients taken coordinatewise. Closure inside the domains
for every finite PSD graph and every defined one-agent elimination is
equivalent to
\begin{equation}
a_i+\frac{u_i\odot a_0}{m}\in\mathscr D_A,\qquad
w_{ij}+\frac{u_i\odot u_j}{m}\in\mathscr D_W,\qquad
m=a_0+\sum_iu_i>0,
\label{eq:cg-kron-star-domain}
\end{equation}
for every finite star inventory of a center anchor \(a_0\) and retained anchors
\(a_i\) in \(\mathscr D_A\), and center weights \(u_i\) and retained weights
\(w_{ij}\) in \(\mathscr D_W\). If only full-space SPD graphs are admitted,
necessity is restricted to inventories producing such graphs. Containment in
\(\mathcal W_H\) does not imply closure inside a smaller family: with the
coefficient set \(\{1\}\) in one dimension, anchor one at three agents and two
unit edges through the middle one, elimination produces the weight \(1/3\) and
retained anchors \(4/3\). \status{ESTABLISHED}

\noindent\emph{Proof.}
The off-diagonal Schur blocks give the first update. In each coordinate,
subtracting the new incident-weight sum from the new diagonal leaves
\begin{equation}
a_i+w_{ik}
-\frac{w_{ik}(w_{ik}+\sum_{j\neq i,k}w_{kj})}{m}
=a_i+\frac{w_{ik}a_k}{m},
\label{eq:cg-kron-scalar-anchor-remainder}
\end{equation}
which is the anchor update. Applying the updates at every star proves
sufficiency of \eqref{eq:cg-kron-star-domain}. Conversely, realize an inventory
as a star with the given retained anchors and weights and eliminate its center;
its pivot is positive, so the outputs must satisfy exactly the displayed
conditions. For several hidden agents, positivity of the hidden principal block
gives positive successive pivots, as in \Cref{prop:cg-normalized-schur-action}.
The one-dimensional example is the update with \(m=3\). \(\square\)

\remarkheading{The shared-kernel version}{cor:cg-kron-shared-kernel}
If \(Q_0\succeq0\) and \(\ker Q_0\subseteq\ker W\) for every
\(W\in\mathscr F\), all forms descend to \(\R^K/\ker Q_0\), the descended anchor
is SPD, and \Cref{thm:cg-independent-kron-classification} applies there with an
ordinary inverse. If \(B_0:\R^{r}\to(\ker Q_0)^\perp\) is an orthonormal
inclusion with \(r=\operatorname{rank}Q_0\), the pair-test property is
equivalent to
\begin{equation}
H_0=B_0(B_0^\top Q_0B_0)^{1/2}O,\qquad
Q_0=H_0H_0^\top,\qquad
W=H_0\operatorname{diag}(d_W)H_0^\top,\quad d_W\geq0,
\label{eq:cg-kron-quotient-cone}
\end{equation}
for one \(O\in O(r)\) and all members \(W\). This is a quotient statement, not
an inverse or a normalized density on the null directions. Without the
shared-kernel hypothesis, \(Q_0=0\) and SPD \(W,W'\) give
\begin{equation}
\bigl[W(t_1W+t_2W')^{-1}W'\bigr]^{-1}=t_1W'^{-1}+t_2W^{-1}\succ0,
\label{eq:cg-kron-zero-anchor-parallel}
\end{equation}
a symmetric output even when \(W\) and \(W'\) do not commute, so the
three-matrix family of \Cref{thm:cg-independent-kron-classification} passes
every test with no common cone. \status{ESTABLISHED}

\noindent\emph{Proof.}
Symmetry and \(Wv=0\) on \(\ker Q_0\) make each form independent of the
representative of a quotient class. Positivity gives
\(\ker(Q_0+t_1W+t_2W')=\ker Q_0\), since a zero of the quadratic annihilates
every nonnegative summand, so the quotient pivot is positive, and whitening in
the coordinates \(B_0\) gives \eqref{eq:cg-kron-quotient-cone}; when \(r=0\) all
forms vanish and the quotient is a point. Expanding the inverse of
\(W(t_1W+t_2W')^{-1}W'\) proves \eqref{eq:cg-kron-zero-anchor-parallel}. A
positive-dimensional null direction has infinite Lebesgue integral, with a
constant or an exponential linear factor along it, so no full-space Gaussian
normalizer is supplied. \(\square\)

\subsection{A graph criterion for unrestricted matrix weights}
\label{sec:cg-kron-graph-boundary}

Now let every present edge weight vary independently over all SPD
\(K\times K\) matrices and restrict the graph instead. Let \(\mathfrak I\) be a
finite undirected multigraph with edge set \(E\) and no self-loops. If anchors
occur, add a ground vertex \(0\), represent each present SPD anchor by an edge
to \(0\), and keep \(0\) among the retained vertices throughout elimination; its
coordinate is set to zero only when the anchored law is formed. In a common
frame,
\begin{equation}
\mathsf E_{\mathrm{aug}}(z)=
\frac12\sum_{e\in E}
(z_{i_e}-z_{j_e})^\top W_e(z_{i_e}-z_{j_e}),\qquad W_e\succ0.
\label{eq:cg-kron-augmented-energy}
\end{equation}
Let \(V_{\mathrm h}\) be the hidden ordinary vertices and \(V_{\mathrm r}\) the
retained vertices, ground included. For each connected component
\(\mathcal Z\) of the graph induced on \(V_{\mathrm h}\), its terminal boundary
is
\begin{equation}
\partial\mathcal Z=\{r\in V_{\mathrm r}:\text{some edge joins }r\text{ to a vertex of }\mathcal Z\}.
\label{eq:cg-kron-terminal-boundary}
\end{equation}

\theoremheading{Universal closure at a fixed retained boundary}{thm:cg-kron-boundary-classification}
Suppose \(\partial\mathcal Z\neq\varnothing\) for every hidden component; this
is exactly the condition that makes the hidden Dirichlet block positive definite
for every SPD assignment, and a component without boundary has a nonzero
constant hidden null vector. For \(K\geq2\), the Schur complement of
\eqref{eq:cg-kron-augmented-energy} is an augmented matrix-weighted Laplacian,
with symmetric nonpositive off-diagonal blocks and diagonal blocks equal to
minus the sum of the off-diagonal blocks, for every independent SPD edge
assignment if and only if
\begin{equation}
|\partial\mathcal Z|\leq2\quad\text{for every hidden component }\mathcal Z,
\label{eq:cg-kron-two-terminal-criterion}
\end{equation}
with ground counted as one terminal. A two-terminal component contributes one
SPD edge and a one-terminal component contributes zero. For \(K=1\) every graph
with nonempty boundaries closes, so the bound is not necessary there.
Sufficiency extends to PSD weights whenever the hidden Dirichlet block remains
SPD, the two-terminal output then being only semidefinite. For a declared
nested collection of hidden sets, universal closure at every declared level is
equivalent to \eqref{eq:cg-kron-two-terminal-criterion} at every such set;
closure at every intermediate one-vertex prefix requires the criterion at every
prefix, because a component with two final terminals can expose three during an
ordering. An excluded graph has some failing SPD assignments, not failure for
every assignment. \status{ESTABLISHED}

\noindent\emph{Proof.}
With retained coordinates zero, zero edge energy forces equality throughout a
hidden component and then zero at its boundary, which proves the
boundary-positivity statement. The hidden Dirichlet block is block diagonal over
the components, and each component's Schur form is the minimum of its energy
with boundary values fixed, hence PSD. Equal boundary values admit the same
constant at every hidden vertex and give zero energy, so a one-terminal
component contributes zero. With two terminals write its symmetric block matrix
as \(M_{\mathcal Z}\); the equations \(M_{11}+M_{12}=0\), \(M_{21}+M_{22}=0\), and
\(M_{21}=M_{12}^\top\) force
\begin{equation}
M_{\mathcal Z}=\begin{pmatrix}W_{\mathcal Z}&-W_{\mathcal Z}\\-W_{\mathcal Z}&W_{\mathcal Z}\end{pmatrix},
\qquad W_{\mathcal Z}=W_{\mathcal Z}^\top\succeq0.
\label{eq:cg-kron-two-terminal-response}
\end{equation}
A nonzero terminal difference with zero energy would have a minimizing hidden
extension with every edge difference zero, and connectedness would force equal
terminal values; hence \(W_{\mathcal Z}\succ0\). Summing the component
contributions and the direct retained edges proves sufficiency, an anchor
increment arising when one terminal is ground.

For necessity, take a hidden center \(4\) with three retained terminals
\(1,2,3\) and legs
\begin{equation}
W_{14}=\begin{pmatrix}1&0\\0&2\end{pmatrix},\qquad
W_{24}=I_2,\qquad
W_{34}=\begin{pmatrix}2&1\\1&2\end{pmatrix}.
\label{eq:cg-kron-three-terminal-weights}
\end{equation}
Its pivot and manufactured coefficient between terminals \(1\) and \(2\) are
\begin{equation}
M=W_{14}+W_{24}+W_{34}=\begin{pmatrix}4&1\\1&5\end{pmatrix},\qquad
W^{\mathrm{Kron}}_{12}=W_{14}M^{-1}W_{24}=
\frac1{19}\begin{pmatrix}5&-1\\-2&8\end{pmatrix},
\label{eq:cg-kron-three-terminal-defect}
\end{equation}
which is not symmetric. If one terminal is ground, give it the \(W_{34}\) leg,
leaving the defect between two ordinary terminals, and direct sums with
\(I_{K-2}\) extend the obstruction to every \(K>2\). Now take any hidden
component with at least three terminals. Select one incident edge from each of
three terminals, join their hidden endpoints by a minimal tree inside the
component, and extend it to a tree containing every hidden vertex of the
component. The three-terminal core has a unique branching vertex, allowing
coincident hidden endpoints and legs of length one, and extra hidden branches
have no terminal, so they minimize to their attachment value and contribute
nothing. On a core leg, minimizing \(\sum_jd_j^\top W_jd_j/2\) subject to
\(\sum_jd_j=d\) gives \(W_jd_j=\lambda\) and
\begin{equation}
W_{\mathrm{leg}}=\Bigl(\sum_jW_j^{-1}\Bigr)^{-1},
\label{eq:cg-kron-flat-series}
\end{equation}
the unique minimum by positivity of the Hessian. Assigning \(lW_{14}\) to each
edge of a leg of length \(l\) realizes \(W_{14}\), and likewise for the other
legs; give the extra tree edges \(I_K\) and temporarily the unused edges zero.
The tree still joins every hidden vertex to a terminal, so its Dirichlet block
is SPD and its reduced core has the defect
\eqref{eq:cg-kron-three-terminal-defect}. Restoring every unused present edge
with weight \(\delta I_K\), the Dirichlet inverse and Schur complement are
continuous at \(\delta=0\), where the determinant is nonzero, so the nonzero
antisymmetric block persists for small \(\delta>0\), with every final weight
SPD. Other hidden components can use scalar identity weights and direct retained
edges are symmetric, so neither cancels the antisymmetric block. For \(K=1\),
closure is \Cref{lem:cg-scalar-kron-closure}.

For the PSD extension, add \(\delta I_K\) to every present weight and let
\(\delta\downarrow0\); continuity of the inverse at the positive limiting
Dirichlet block preserves the Schur matrix, and symmetry, zero row sums, and
PSD weights are closed conditions. Direct and successive elimination agree by
\Cref{prop:cg-normalized-schur-action}, which gives the nested statement, and
the prefix statement is the same criterion applied at each prefix. \(\square\)

Flatness here means plain differences in a common frame. A graph coboundary can
first be trivialized; genuinely nonflat links are the subject of
\Cref{sec:cg-nonflat-series-kron}.

\corollaryheading{Marginalizing prior-carrying agents}{cor:cg-kron-anchored-agents}
Let \(K\geq2\), and suppose every hidden agent carries a present SPD anchor,
for instance a proper agent-level prior self term. Then the Schur complement of
\eqref{eq:cg-kron-augmented-energy} is an augmented matrix-weighted Laplacian
for every independent SPD assignment of the edge weights and anchors if and only
if every hidden component is adjacent to at most one retained ordinary agent.
Such a component adds an SPD increment to the anchor of that agent, or nothing
when it has no retained ordinary neighbor. On every other graph, in particular
when a hidden agent relays between two retained agents, the failing assignments
contain an open dense subset of full Lebesgue measure in the SPD parameter
domain. \status{ESTABLISHED}

\noindent\emph{Proof.}
Each present anchor is an edge to ground, so ground lies in
\(\partial\mathcal Z\) for every hidden component, and
\eqref{eq:cg-kron-two-terminal-criterion} holds exactly when \(\mathcal Z\) has
at most one retained ordinary neighbor. By
\Cref{thm:cg-kron-boundary-classification}, a component whose terminals are one
retained agent \(r\) and ground contributes one SPD edge from \(r\) to ground,
an anchor increment at \(r\), and a component whose only terminal is ground
contributes zero. On an excluded graph the theorem supplies an SPD assignment
at which some retained off-diagonal block is nonsymmetric. By the adjugate
formula every entry of the antisymmetric part of that block is a polynomial in
the weight entries divided by \(\det\Lambda_{\mathrm{hh}}>0\); the SPD parameter
domain is an open connected subset of a Euclidean space, and a polynomial that
is nonzero at one of its points vanishes there only on a closed nowhere-dense
set of Lebesgue measure zero. For a single relay with hidden anchor \(Q_0\) and
incident weights \(W,W'\), \eqref{eq:cg-exact-pair-antisymmetry} with
\(t_1=t_2=1\) shows that the manufactured block is symmetric exactly when
\(WQ_0^{-1}W'=W'Q_0^{-1}W\). \Cref{prop:cg-kron-leaves-family} is such a relay
with \(Q_0=I_2\): it is the star \eqref{eq:cg-kron-three-terminal-weights} with
its \(W_{24}=I_2\) leg grounded, and its manufactured coefficient is
\(W_{14}M^{-1}W_{34}\). \(\square\)

The interaction family only requires \(A_i\succeq0\), and a directed forest
places proper priors at its roots alone (\Cref{prop:gauss-gain-absorption}), so
the hypothesis is a declared restriction. When every agent does carry an SPD
self term, marginalizing an agent that relays between two retained agents
generically leaves the positive matrix-Laplacian family, while marginalizing a
hidden component attached to a single retained agent stays in it and only
raises that agent's anchor.

\subsection{A proper Gaussian beyond every positive two-node decomposition}
\label{sec:cg-kron-positive-pair-obstruction}

The three-terminal star also separates failure of the difference-weight form
from failure of every positive pair representation on the same retained nodes.
\citet[Sec.~4.2]{HansenGhrist2019} show that Kron reduction is impossible for
general cellular sheaves: their witness is a star with three boundary vertices,
one-dimensional boundary stalks, and a two-dimensional internal stalk, whose
boundary section space no sheaf on the boundary vertices alone reproduces. For
one-dimensional vertex stalks, \citet[Thm.~4.3]{HansenGhrist2019} prove the opposite: matrices
of factor width at most two, which are exactly the Laplacians of such sheaves
\citep[Sec.~3.2.2]{HansenGhrist2019}, are closed under Schur complements. The
positive two-node class below is the block analog of factor width two, with
\(K\)-dimensional node blocks, and its dual cone consists of the symmetric
matrices whose principal submatrices on at most two nodes are positive
semidefinite. Relative to that work, the proposition supplies a proper SPD
Gaussian law rather than the section space of a singular sheaf Laplacian, with
two-dimensional node blocks and an exact rational certificate in the dual cone.

\propositionheading{Rational obstruction to positive pair decomposition}{prop:cg-kron-positive-pair-obstruction}
There is a proper four-agent Gaussian with SPD edge weights and retained anchors
whose marginal on three two-dimensional nodes is not a sum of PSD quadratic
terms each supported on at most two of those nodes. Hence new invertible links,
rectangular pair residuals, parallel edges, self-loops, and freely
redistributed PSD anchors cannot represent that marginal within the positive
two-node class, and neither can block-diagonal invertible changes of the
retained frames, which preserve the class by congruence. The marginal action is
of M\"obius order two: every M\"obius potential
\eqref{eq:rg-mobius-potentials} of order three or more of a Gaussian marginal
action vanishes, so the obstruction concerns positivity of the pair terms, not
interaction order. \status{ESTABLISHED}

\noindent\emph{Proof.}
Let \(L_{\mathrm{st}}\) be the six-dimensional retained Schur Laplacian of the
unanchored star \eqref{eq:cg-kron-three-terminal-weights}, and put
\begin{equation}
\mathsf J=\begin{pmatrix}0&1\\-1&0\end{pmatrix},\qquad
t=\frac1{104},\qquad \delta=\frac1{20000}.
\label{eq:cg-kron-rational-separator-constants}
\end{equation}
Define a symmetric matrix \(Y\) in three two-dimensional blocks by
\begin{equation}
\begin{aligned}
Y_{ii}&=(1+t^2)I_2,\\
Y_{12}&=I_2+t\mathsf J,\qquad Y_{21}=I_2-t\mathsf J,\\
Y_{ij}&=I_2\quad\text{for the other distinct pairs}.
\end{aligned}
\label{eq:cg-kron-rational-separator}
\end{equation}
Every principal submatrix of \(Y\) on one or two nodes is PSD: for diagonal
blocks \(\alpha I\) and a cross block \(C\) the quadratic is at least
\(\alpha(\|u\|^2+\|v\|^2)-2\|C\|\|u\|\|v\|\geq0\) when \(\|C\|\leq\alpha\), and
\((I_2+t\mathsf J)^\top(I_2+t\mathsf J)=(1+t^2)I_2\) gives the special cross norm
\(\sqrt{1+t^2}\leq1+t^2\), the others being one. Consequently
\(\Tr(YP)\geq0\) for every PSD matrix \(P\) supported on at most two nodes,
since \(\Tr(YP)=\Tr(Y_{TT}P_{TT})\) for \(P\) supported on a node set \(T\): this
is the dual-cone membership of \(Y\), and every sum of positive one-node or
two-node terms has nonnegative pairing with \(Y\).

Substitution into \(W_{i4}-W_{i4}M^{-1}W_{i4}\) gives diagonal traces
\(36/19\), \(29/19\), \(39/19\), hence \(\Tr L_{\mathrm{st}}=104/19\). Write
\(Y=(\mathbf1\mathbf1^\top)\otimes I_2+tY_1+t^2I_6\), with
\((Y_1)_{12}=\mathsf J\), \((Y_1)_{21}=-\mathsf J\), and other blocks zero. The
constant-vector kernel of \(L_{\mathrm{st}}\) makes the first pairing vanish,
and since
\(W^{\mathrm{Kron}}_{12}-(W^{\mathrm{Kron}}_{12})^\top=\mathsf J/19\),
\begin{equation}
\Tr(Y_1L_{\mathrm{st}})
=\Tr\bigl(\mathsf J(W^{\mathrm{Kron}}_{12}-(W^{\mathrm{Kron}}_{12})^\top)\bigr)=-\frac2{19},
\qquad
\Tr(YL_{\mathrm{st}})=\frac{-2t+104t^2}{19}=-\frac1{1976}.
\label{eq:cg-kron-negative-dual-pairing}
\end{equation}
Put the anchor \(\delta I_2\) at each of the three retained vertices of the
star, keeping its hidden anchor zero. Zero total energy forces the anchored
terminal coordinates and then the connected hidden coordinate to vanish, so the
full precision is SPD, and its retained precision is
\(\Lambda_{\mathrm r}=L_{\mathrm{st}}+\delta I_6\succ0\), but
\begin{equation}
\Tr(Y\Lambda_{\mathrm r})<-\frac1{1976}+\frac7{20000}<0,
\qquad
\Tr Y=6(1+t^2)<7,
\label{eq:cg-kron-proper-dual-obstruction}
\end{equation}
the last step by \(7\cdot1976=13832<20000\). This contradicts the nonnegative
pairing of every positive two-node decomposition. All constants are rational,
and \eqref{eq:cg-kron-normalized-prior} supplies the finite normalizer. A
block-diagonal invertible change of retained frames maps positive one-node and
two-node terms to terms of the same kind. For the M\"obius statement, a Gaussian
action is a sum of functions of at most two node blocks; for a function of the
blocks in a set \(D\) and a block \(a\notin D\), the M\"obius potential of any
set containing \(a\) vanishes, because pairing each subset \(B\) without \(a\)
with \(B\cup\{a\}\) gives equal terms of opposite sign. \(\square\)

Mixing distinct retained nodes, adding latent vertices, using signed pair terms,
or allowing positive terms on three or more nodes changes the representation
problem, and the SPD marginal remains a valid complete quadratic Gaussian
action. The hyperedges that exact closure generically requires for general laws
(\Cref{sec:rg-hypergraph-closure}) do not arise for Gaussian marginals.

\subsection{Nonflat series and cycle marginalization}
\label{sec:cg-nonflat-series-kron}

We now retain link data instead of requiring a flat difference presentation.
Loop holonomy at \(c_*\) is in scope here, and the content is Gaussian.

\definitionheading{Marked nonflat series class}{def:cg-nonflat-series-class}
Fix a finite multigraph on \(V\), one orientation of each edge copy
\(e=(i,j)\), invertible links \(\Theta_e\in\GL(K)\) mapping the fiber at \(j\) to
the fiber at \(i\), SPD weights \(W_e\), PSD anchors \(A_i\), sources
\(\mathsf h_i\), and a scalar offset \(\mathfrak c\), with action
\begin{equation}
\mathsf E(z)=\mathfrak c
+\frac12\sum_{e=(i,j)}
(z_i-\Theta_ez_j)^\top W_e(z_i-\Theta_ez_j)
+\frac12\sum_i z_i^\top A_i z_i-\sum_i \mathsf h_i^\top z_i .
\label{eq:cg-nonflat-series-action}
\end{equation}
A split \(V=V_{\mathrm r}\sqcup V_{\mathrm h}\) belongs to the marked nonflat
series class when every hidden vertex has total degree two, a self-loop counted
twice, and zero anchor, the retained anchors are kept throughout elimination,
and the assembled precision \(\Lambda\) is SPD. Hidden sources need not vanish.

\lemmaheading{Properness of a marked nonflat graph}{lem:cg-nonflat-series-properness}
In each connected component choose a root \(r\) and a spanning tree, with tree
transports \(\Theta_{i\leftarrow r}\) and \(\Theta_{r\leftarrow r}=I\). Then
\begin{equation}
\Lambda\succ0
\quad\Longleftrightarrow\quad
\bigcap_{\substack{e=(i,j)\\\text{non-tree}}}
\ker(\Theta_{i\leftarrow r}-\Theta_e\Theta_{j\leftarrow r})
\ \cap\ \bigcap_i\ker(A_i\Theta_{i\leftarrow r})=\{0\}
\quad\text{in each component}.
\label{eq:cg-nonflat-series-properness}
\end{equation}
One retained SPD anchor in every component suffices. An all-hidden unanchored
cycle with holonomy \(\mathrm{Hol}\) is admitted exactly when \(I-\mathrm{Hol}\)
is invertible, and \(\mathrm{Hol}\neq I\) alone does not suffice.
\status{ESTABLISHED}

\noindent\emph{Proof.}
The quadratic form of \(\Lambda\) is a sum of nonnegative residual and anchor
terms, so its kernel consists of the vectors with \(z_i=\Theta_ez_j\) on every
edge and \(A_iz_i=0\) at every vertex. The tree edges force
\(z_i=\Theta_{i\leftarrow r}u\) with \(u=z_r\), and the remaining conditions are
the displayed intersection. An SPD anchor at a vertex \(i\) forces
\(\Theta_{i\leftarrow r}u=0\), hence \(u=0\). On a cycle with one non-tree edge
the condition reads \((I-\mathrm{Hol})u=0\) after conjugation by an invertible
tree transport, so it has only the zero solution exactly when \(I-\mathrm{Hol}\)
is invertible; \(\mathrm{Hol}=\operatorname{diag}(1,2)\) is a nonidentity
holonomy with a nonzero kernel. \(\square\)

\propositionheading{Exact nonflat series pivot}{prop:cg-nonflat-series-pivot}
Let \(i_0\) be a hidden vertex of the class of
\Cref{def:cg-nonflat-series-class}, and orient its two non-loop edges so that
their residuals are \(x-\Theta_1z_{i_0}\) and \(z_{i_0}-\Theta_2y\), with SPD
weights \(W_1,W_2\), where \(x=z_{i_1}\) and \(y=z_{i_2}\) may be coordinates of
the same vertex. Set
\begin{equation}
M=\Theta_1^\top W_1\Theta_1+W_2,\qquad
W_{\mathrm{new}}=(W_1^{-1}+\Theta_1W_2^{-1}\Theta_1^\top)^{-1},\qquad
\Theta_{\mathrm{new}}=\Theta_1\Theta_2 .
\label{eq:cg-nonflat-series-weight}
\end{equation}
Eliminating \(z_{i_0}\) replaces the two edge terms by
\(\frac12(x-\Theta_{\mathrm{new}}y)^\top W_{\mathrm{new}}(x-\Theta_{\mathrm{new}}y)\),
and the source and scalar updates are
\begin{equation}
\begin{aligned}
\mathsf h_{i_1}^{\mathrm{new}}&=\mathsf h_{i_1}+W_1\Theta_1M^{-1}\mathsf h_{i_0},\\
\mathsf h_{i_2}^{\mathrm{new}}&=\mathsf h_{i_2}+\Theta_2^\top W_2M^{-1}\mathsf h_{i_0},\\
\mathfrak c^{\mathrm{new}}
&=\mathfrak c+\frac12\log\det M-\frac K2\log(2\pi)
                         -\frac12\mathsf h_{i_0}^\top M^{-1}\mathsf h_{i_0}.
\end{aligned}
\label{eq:cg-nonflat-series-source-scalar}
\end{equation}
If \(i_1=i_2\), both source increments add at that vertex and the new edge is
the marked self-loop \((\Theta_{\mathrm{new}},W_{\mathrm{new}})\). No new
anchor is generated. \status{ESTABLISHED}

\noindent\emph{Proof.}
The hidden terms are
\begin{equation}
\frac12z_{i_0}^\top Mz_{i_0}-z_{i_0}^\top q
+\frac12x^\top W_1x+\frac12y^\top\Theta_2^\top W_2\Theta_2y,
\qquad q=\Theta_1^\top W_1x+W_2\Theta_2y+\mathsf h_{i_0}.
\label{eq:cg-nonflat-series-square-data}
\end{equation}
Since \(M\succ0\), completing the square multiplies the remaining density by
\[
(2\pi)^{K/2}(\det M)^{-1/2}\exp\bigl(q^\top M^{-1}q/2\bigr).
\]
Let
\(\bar W=W_1-W_1\Theta_1M^{-1}\Theta_1^\top W_1\). From
\(M=\Theta_1^\top W_1\Theta_1+W_2\) one has
\(\bar W\Theta_1=W_1\Theta_1M^{-1}W_2\), and consequently
\begin{equation}
\begin{aligned}
\bar W(W_1^{-1}+\Theta_1W_2^{-1}\Theta_1^\top)&=I,\\
W_2-W_2M^{-1}W_2&=\Theta_1^\top \bar W\Theta_1 .
\end{aligned}
\label{eq:cg-nonflat-series-matrix-identities}
\end{equation}
The first identity gives \(\bar W=W_{\mathrm{new}}\); the second and
\(\bar W\Theta_1=W_1\Theta_1M^{-1}W_2\) identify the endpoint diagonal and cross
terms as the new edge energy, and the linear and constant terms of
\(q^\top M^{-1}q/2\) give \eqref{eq:cg-nonflat-series-source-scalar}.
Substituting \(x=y\) gives the loop. Reversing an edge uses
\((\Theta,W)\mapsto(\Theta^{-1},\Theta^\top W\Theta)\), because the old
residual is minus \(\Theta\) times the reversed one, so the displayed
orientation costs nothing. \(\square\)

\theoremheading{Arbitrary-order nonflat series closure}{thm:cg-nonflat-series-kron}
In the class of \Cref{def:cg-nonflat-series-class}, eliminating any subset of
hidden vertices in any order gives the exact normalized Gaussian marginal within
the same class: every surviving hidden vertex still has degree two and zero
anchor, the remaining precision is SPD, and distinct parallel paths remain
distinct edges. Along an oriented path \(v_0,\ldots,v_m\), with \(\Theta_i\)
mapping from \(v_i\) to \(v_{i-1}\), put \(\Theta_{\leq0}=I\) and
\(\Theta_{\leq j}=\Theta_1\cdots\Theta_j\); the fused data are
\begin{equation}
\Theta_{\mathrm{path}}=\Theta_{\leq m},\qquad
W_{\mathrm{path}}^{-1}
=\sum_{i=1}^m\Theta_{\leq i-1}W_i^{-1}\Theta_{\leq i-1}^\top .
\label{eq:cg-nonflat-path-formula}
\end{equation}
A returning path has retained loop holonomy \(\mathrm{Hol}=\Theta_{\leq m}\)
and quadratic
\begin{equation}
Q_{\mathrm{loop}}=(I-\mathrm{Hol})^\top W_{\mathrm{path}}(I-\mathrm{Hol}).
\label{eq:cg-nonflat-retained-loop}
\end{equation}
A component meeting \(V_{\mathrm r}\) decomposes into maximal hidden paths
between retained endpoints, which may coincide; a cycle with one retained vertex
becomes one loop, and with several retained vertices its ordered coarse link
product equals its original cycle product, so each cycle meeting the retained
set keeps its holonomy mark. An all-hidden cycle integrates away exactly under
the condition of \Cref{lem:cg-nonflat-series-properness}; its raw partition
factor survives, but its normalized marginal on the empty coordinate set does
not identify its holonomy, and a retained \(Q_{\mathrm{loop}}\) need not
identify the separate marks \(\mathrm{Hol}\) and \(W_{\mathrm{path}}\). These
marks are structural data that the Gaussian law does not determine.
\status{ESTABLISHED}

\noindent\emph{Proof.}
A degree-two hidden vertex without a loop has two incident edges counted with
multiplicity, so \Cref{prop:cg-nonflat-series-pivot} applies. If its neighbors
differ, one edge replaces two and degrees are preserved; if they coincide, a
loop carries the same two incidences. No hidden anchor is created. A degree-two
vertex already carrying a loop is an isolated component, arising initially or as
the last vertex of an all-hidden cycle; its loop block is SPD by global
positivity and integrates as one Gaussian block, contributing only the scalar
of \eqref{eq:cg-kron-complete-schur}. Positivity of every remaining Schur
complement holds by \Cref{prop:cg-normalized-schur-action}, which proves closure
by induction for every order. Write a path segment as \((\Theta,C)\) with
\(C=W^{-1}\); the series operation is
\begin{equation}
(\Theta,C)*(\Theta',C')
=(\Theta\Theta',\ C+\Theta C'\Theta^\top).
\label{eq:cg-nonflat-series-composition}
\end{equation}
For a third pair \((\Theta'',C'')\) either parenthesization gives transport
\(\Theta\Theta'\Theta''\) and covariance
\(C+\Theta C'\Theta^\top+\Theta\Theta'C''\Theta'^\top\Theta^\top\), which proves
associativity and, by induction, \eqref{eq:cg-nonflat-path-formula}; each
summand is SPD. Equal endpoints give \eqref{eq:cg-nonflat-retained-loop}. The
scalar and source terms agree for all orders by the Gaussian integration
argument of \Cref{prop:cg-normalized-schur-action}. The coarse link product of a
cycle is the ordered product of the fused links, which is the original cycle
product. Different pairs \((\mathrm{Hol},W_{\mathrm{path}})\) can give the same
\(Q_{\mathrm{loop}}\), for instance \((2I,W)\) and \((3I,W/4)\), both with
\(Q_{\mathrm{loop}}=W\). \(\square\)

\propositionheading{Gauge and rescaling covariance with scalar factors}{prop:cg-nonflat-series-gauge}
For nodewise passive coordinates \(z_i'=g_iz_i\), \(g_i\in\GL(K)\), the data
transform as
\begin{equation}
\begin{aligned}
\Theta_e'&=g_i\Theta_eg_j^{-1},&
W_e'&=g_i^{-\top}W_eg_i^{-1},\\
A_i'&=g_i^{-\top}A_ig_i^{-1},&
\mathsf h_i'&=g_i^{-\top}\mathsf h_i,
\end{aligned}
\label{eq:cg-nonflat-series-gauge-law}
\end{equation}
and the series rule commutes with the change:
\(M'=g_{i_0}^{-\top}Mg_{i_0}^{-1}\),
\(W_{\mathrm{new}}'=g_{i_1}^{-\top}W_{\mathrm{new}}g_{i_1}^{-1}\), and
\(\Theta_{\mathrm{new}}'=g_{i_1}\Theta_{\mathrm{new}}g_{i_2}^{-1}\). With the raw
scalar \(\mathfrak c\) unchanged,
\begin{equation}
\begin{aligned}
\mathfrak c_{\mathrm r}'
 &=\mathfrak c_{\mathrm r}-\sum_{i\in V_{\mathrm h}}\log|\det g_i|,\qquad
Z'=\Bigl(\prod_{i\in V}|\det g_i|\Bigr)Z,\\
p_{\mathrm r}'(x')&=|\det g_{\mathrm r}|^{-1}p_{\mathrm r}(g_{\mathrm r}^{-1}x'),
\qquad g_{\mathrm r}=\operatorname{blkdiag}_{i\in V_{\mathrm r}}g_i ;
\end{aligned}
\label{eq:cg-nonflat-series-jacobians}
\end{equation}
if instead the raw measure itself is transformed,
\(\mathfrak c'=\mathfrak c+\sum_{i\in V}\log|\det g_i|\), then \(Z'=Z\) and
\(\mathfrak c_{\mathrm r}'=\mathfrak c_{\mathrm r}+\sum_{i\in V_{\mathrm r}}\log|\det g_i|\),
and both conventions give the same normalized law. A retained loop transforms
by \(\mathrm{Hol}'=g_r\mathrm{Hol}g_r^{-1}\). An invertible retained rescaling
\(x'=gx\) has the normalized pushforward
\begin{equation}
\mu_{\mathrm r}'=g\mu_{\mathrm r},\qquad \Sigma_{\mathrm r}'=g\Sigma_{\mathrm r}g^\top,\qquad
p_{\mathrm r}'(x')=|\det g|^{-1}p_{\mathrm r}(g^{-1}x'),
\label{eq:cg-nonflat-series-rescaling}
\end{equation}
which preserves the marked graph presentation when \(g\) is block diagonal and
may mix vertices otherwise. \status{ESTABLISHED}

\noindent\emph{Proof.}
Substitution into each residual proves \eqref{eq:cg-nonflat-series-gauge-law}.
The inverse covariance of \eqref{eq:cg-nonflat-series-weight} transforms by
\begin{equation}
W_1'^{-1}+\Theta_1'W_2'^{-1}\Theta_1'^\top
=g_{i_1}(W_1^{-1}+\Theta_1W_2^{-1}\Theta_1^\top)g_{i_1}^\top .
\label{eq:cg-nonflat-series-covariance-gauge}
\end{equation}
The source increments transform as endpoint covectors,
\(\mathsf h_{i_0}^\top M^{-1}\mathsf h_{i_0}\) is invariant, and
\(\det M'=\det M/|\det g_{i_0}|^2\), which gives the hidden scalar shift.
Changing variables in the full integral gives the total Jacobian and then the
retained one; adding the total log-Jacobian to the initial action gives the
raw-measure convention. The loop and rescaling laws are the same change of
variables. \(\square\)

A restriction of the links to a matrix subgroup must be closed under products
and inverses, and gauges in that subgroup preserve the restricted domain. No
preferred physical scale transformation is selected by this finite pushforward.

\propositionheading{A noncommuting nonflat proper example}{prop:cg-nonflat-series-example}
Take one retained vertex \(1\), hidden vertices \(2,3\), and two distinct
two-edge cycles \(1-2-1\) and \(1-3-1\), whose edge copies are distinct
interactions. Give vertex \(1\) the anchor \(I_2\), the hidden vertices zero
anchors, and put
\begin{equation}
\Theta_{\mathrm a}=\begin{pmatrix}1&1\\0&1\end{pmatrix},\qquad
\Theta_{\mathrm b}=\begin{pmatrix}1&0\\1&1\end{pmatrix},\qquad
W_{\mathrm a}=\begin{pmatrix}1&0\\0&2\end{pmatrix},\qquad
W_{\mathrm b}=\begin{pmatrix}2&1\\1&2\end{pmatrix}.
\label{eq:cg-nonflat-series-example-data}
\end{equation}
The first cycle has links \(\Theta_{\mathrm a},I_2\) and weights
\(W_{\mathrm a},W_{\mathrm b}\); the second has links \(\Theta_{\mathrm b},I_2\)
and weights \(I_2,I_2\). With \(\mathsf h_2=(1,0)^\top\), all other sources
zero, and \(\mathfrak c=0\), this is a proper Gaussian with noncommuting loop
holonomies and coefficients outside every cone \(\mathcal W_H\) containing its
anchor, and its exact marginal on vertex \(1\) is given by
\eqref{eq:cg-nonflat-series-example-law}. \status{ESTABLISHED}

\noindent\emph{Proof.}
Every weight is SPD and the retained SPD anchor gives
\eqref{eq:cg-nonflat-series-properness}. The holonomies satisfy
\(\Theta_{\mathrm a}\Theta_{\mathrm b}-\Theta_{\mathrm b}\Theta_{\mathrm a}
=\operatorname{diag}(1,-1)\neq0\), so no based conjugation makes both flat or
diagonal. The triple \(I_2,W_{\mathrm a},W_{\mathrm b}\) lies in no common cone:
containing \(I_2\) forces a common orthogonal basis, but
\(W_{\mathrm a}W_{\mathrm b}-W_{\mathrm b}W_{\mathrm a}
=\begin{psmallmatrix}0&-1\\1&0\end{psmallmatrix}\). This fixed-chart
coefficient statement and the gauge-invariant holonomy statement are distinct.
The hidden pivots and loop weights are
\begin{equation}
\begin{aligned}
M_2&=\begin{pmatrix}3&2\\2&5\end{pmatrix},&
\widehat W_2&=\frac1{11}\begin{pmatrix}7&-2\\-2&10\end{pmatrix},\\
M_3&=\begin{pmatrix}3&1\\1&2\end{pmatrix},&
\widehat W_3&=\frac15\begin{pmatrix}3&-1\\-1&2\end{pmatrix},
\end{aligned}
\label{eq:cg-nonflat-series-example-pivots}
\end{equation}
with \(\det M_2=11\), \(\det M_3=5\), and loop quadratics
\(\operatorname{diag}(0,7/11)\) and \(\operatorname{diag}(2/5,0)\). The formulas
above give
\begin{equation}
\begin{aligned}
\Lambda^{\mathrm{eff}}&=\operatorname{diag}(7/5,18/11),&
\mathsf h^{\mathrm{eff}}&=(1,-3/11)^\top,\\
\mathfrak c^{\mathrm{eff}}
 &=\frac12\log55-2\log(2\pi)-\frac5{22},&
\mu_1&=(5/7,-1/6)^\top,\\
\Sigma_1&=\operatorname{diag}(5/7,11/18),&
Z&=(2\pi)^3\exp(17/28)/\sqrt{126}.
\end{aligned}
\label{eq:cg-nonflat-series-example-law}
\end{equation}
Indeed \(\det\Lambda=11\cdot5\cdot(126/55)=126\),
\(M_2^{-1}\mathsf h_2=(5,-2)^\top/11\), and the completed-square value is
\(5/11+5/7+1/22=17/14\), which verifies source, covariance, scalar, and
normalization in the same proper law. \(\square\)

The zero hidden anchor is what makes the series rule produce no new anchor.
Under admitted positive-semidefinite self terms, anchored hidden vertices still
close, with a presentation that is no longer unique.

\propositionheading{Anchored hidden vertices}{prop:cg-nonflat-series-anchored}
Keep the setting of \Cref{prop:cg-nonflat-series-pivot} with distinct
endpoints, but give the hidden vertex a PSD anchor \(A_{i_0}\), and put
\(M=\Theta_1^\top W_1\Theta_1+W_2+A_{i_0}\). Eliminating \(z_{i_0}\) replaces the
two edge terms and the anchor by
\(\frac12x^\top \bar Wx-x^\top Ry+\frac12y^\top Qy\), with
\[
\bar W=W_1-W_1\Theta_1M^{-1}\Theta_1^\top W_1\succ0,\qquad
R=W_1\Theta_1M^{-1}W_2\Theta_2,\qquad
Q=\Theta_2^\top(W_2-W_2M^{-1}W_2)\Theta_2,
\]
where \(R\) is invertible and \(Q-R^\top \bar W^{-1}R\succeq0\), so that
\[
\tfrac12(x-\bar W^{-1}Ry)^\top \bar W(x-\bar W^{-1}Ry)+\tfrac12y^\top(Q-R^\top \bar W^{-1}R)y
\]
is one invertible link with SPD weight plus a generated PSD anchor at
\(i_2\); the source and scalar updates are those of
\eqref{eq:cg-nonflat-series-source-scalar} with this \(M\). For
\(A_{i_0}=0\) this is the series rule, with \(\bar W=W_{\mathrm{new}}\),
\(\bar W^{-1}R=\Theta_{\mathrm{new}}\), and zero anchor. For \(A_{i_0}\neq0\) new
anchors arise: in the scalar chain with unit edges, identity links, and hidden
anchor \(\alpha>0\), the pivot is \(M=2+\alpha\) and the retained quadratic
matrix is
\begin{equation}
\begin{pmatrix}1-1/M&-1/M\\-1/M&1-1/M\end{pmatrix}
=\frac1M\begin{pmatrix}1&-1\\-1&1\end{pmatrix}
+\frac{\alpha}{M}I_2 .
\label{eq:cg-nonflat-series-hidden-anchor-boundary}
\end{equation}
The presentation is not unique: the flat link with weight \(1/M\) and anchors
\(\alpha/M\) at both endpoints, and the link \(1/(1+\alpha)\) with weight
\((1+\alpha)/M\) and anchor \(\alpha/(1+\alpha)\) at \(i_2\), both reproduce
\eqref{eq:cg-nonflat-series-hidden-anchor-boundary}. \status{ESTABLISHED}

\noindent\emph{Proof.}
The hidden terms are those of \eqref{eq:cg-nonflat-series-square-data} with the
anchor added to \(M\), and completing the square leaves
\(\frac12x^\top W_1x+\frac12y^\top\Theta_2^\top W_2\Theta_2y-\frac12q^\top M^{-1}q\),
whose blocks are \(\bar W\), \(-R\), and \(Q\); the source and scalar terms are as
before. The matrix
\(\begin{psmallmatrix}W_1&W_1\Theta_1\\\Theta_1^\top W_1&M\end{psmallmatrix}\)
is SPD because its Schur complement at \(W_1\) is \(W_2+A_{i_0}\succ0\), so its
Schur complement \(\bar W\) at \(M\) is SPD. \(R\) is a product of invertible
matrices. The eliminated quadratic is the minimum over \(z_{i_0}\) of a
nonnegative form, hence PSD, and with \(\bar W\succ0\) its Schur complement
\(Q-R^\top \bar W^{-1}R\) is PSD; completing the square in \(x\) gives the displayed
presentation. For \(A_{i_0}=0\), \(\bar W=W_{\mathrm{new}}\) by
\eqref{eq:cg-nonflat-series-matrix-identities}, \(R=W_{\mathrm{new}}\Theta_1\Theta_2\)
by \(\bar W\Theta_1=W_1\Theta_1M^{-1}W_2\), and \(Q=\Theta_2^\top\Theta_1^\top
W_{\mathrm{new}}\Theta_1\Theta_2=R^\top \bar W^{-1}R\). The scalar chain is direct
substitution: \(\bar W=Q=1-1/M\) and \(R=1/M\), and both presentations expand to
\eqref{eq:cg-nonflat-series-hidden-anchor-boundary}. \(\square\)

In the flat graph theorem a hidden anchor is an edge to ground, which can
create a third terminal, and in the fixed-anchor theorem the positive hidden
anchor is the premise that forces the commutation constraint. Degree two in the
nonflat theorem guarantees every intermediate order, whereas the graph criterion
concerns each selected final hidden set; these quantifiers are distinct. The
nonflat series closure assumes zero hidden anchors and does not cover
prior-carrying hidden agents, for which the presentation of
\Cref{prop:cg-nonflat-series-anchored} carries a link that is not determined by
the marginal.
